\scuAbstractZhHeading
高分子链在不同尺度上的构象决定了材料结构与宏观响应之间的联系。链段数增加时，单链的空间尺寸并不简单地按轮廓长度线性增长，而是受到局部相关性、排除体积以及溶剂环境等因素的共同影响。建立清楚的尺度关系，有助于区分由链结构决定的规律与由实验条件引起的差异。本文以理想链为起点，讨论均方末端距、回转半径与有效链段长度之间的关系，并将其与良溶剂条件下的尺度描述进行比较。

研究首先明确自由连接链的基本假设，将链的末端矢量写为各链段矢量之和。通过展开矢量内积，说明独立且各向同性的链段为何使交叉项的统计平均为零，从而得到均方末端距与链段数成正比的结果。这里的独立性是模型前提，不能直接推广到存在短程刚性或长程相互作用的真实链。对于具有局部取向关联的链，本文采用有效链段的思想重新解释模型参数，强调统计链段长度并不等于化学重复单元的几何长度。

在此基础上，本文讨论理想链与自回避链的差别。理想链允许不同链段在空间中相互穿越，因此能够给出简洁的解析关系；真实链的排除体积则使可访问构象发生变化。在足够长的链和适当溶剂条件下，链尺寸可用幂律形式描述，但指数的物理含义必须与空间维度、溶剂质量和观测尺度一并解释。单独拟合一条双对数直线，不能证明体系已进入渐近尺度区间，也不能据此排除有限链长和链刚性的影响。

为展示这些区别，本文构造一组明确标注为演示用途的无量纲模型曲线，并列出理想链与良溶剂链在假设、尺度指数及适用范围方面的差异。曲线来自指定公式，不代表实测结果，也不承担材料性能预测的作用。示例同时展示公式、图、表、脚注与分级标题的组织方式，使模型推导与文档排版能够分别检查。对可能的实验验证，本文提出在独立测定分子量分布与溶剂条件的基础上，比较不同链长区间的尺寸变化，并检验拟合指数对数据范围的敏感性。

分析表明，理想链关系的价值在于提供可解释的参照，而非为所有高分子体系提供统一答案。将模型假设、统计量定义和数据来源写清楚，可以减少参数拟合中的过度解释。真实材料研究还需要考虑多分散性、链间作用与测量误差。本文属于理论与排版示例，不包含原创实验数据；实际学位论文应以作者完成的研究替换相关内容，并依据学院要求补足论证、文献与实验依据。

\par\vspace{20bp}\noindent{\scuHei 关键词：}\scuKeywordsZh
